% !TeX program = lualatex
\documentclass[12pt,a4paper]{article}
\usepackage[utf8]{inputenc}
\usepackage[T1]{fontenc}
\usepackage{lmodern}
\usepackage[ngerman]{babel}
\usepackage{amsmath,amssymb,amsthm,mathtools}
\usepackage{geometry}
\usepackage{booktabs}
\usepackage{hyperref}
\usepackage{url}
\usepackage{tabularx}
\usepackage{array}
\usepackage{makecell}
\usepackage{ragged2e}
\usepackage{bm}
\usepackage{slashed}
\usepackage{tikz}
\usetikzlibrary{arrows.meta, positioning, calc, shapes.geometric}
\usepackage{pgfplots}
\pgfplotsset{compat=1.18}
\usepackage{graphicx}
\usepackage{caption}
\usepackage{subcaption}
\usepackage{float}
\usepackage{nicefrac}

% Theorem-Umgebungen (deutsch)
\newtheorem{theorem}{Satz}[section]
\newtheorem{lemma}[theorem]{Lemma}
\newtheorem{definition}[theorem]{Definition}
\newtheorem{corollary}[theorem]{Korollar}
\newtheorem{proposition}[theorem]{Proposition}
\theoremstyle{remark}
\newtheorem{remark}[theorem]{Bemerkung}
\newtheorem{example}[theorem]{Beispiel}

% Geometrie
\geometry{a4paper, margin=1in}

% YUCT-Makros (unverändert)
\newcommand{\Keff}{K_{\mathrm{eff}}}
\newcommand{\Keffzero}{K_{\mathrm{eff},0}}
\newcommand{\kappac}{\kappa_c}
\newcommand{\eps}{\varepsilon}
\newcommand{\df}{d_{\!f}}
\newcommand{\constmin}{C_{\min}}
\newcommand{\dint}{\ensuremath{D\!+\!I\!\cdot\!R}}
\newcommand{\Mpl}{M_{\mathrm{Pl}}}
\newcommand{\mpl}{m_{\mathrm{Pl}}}
\newcommand{\Lpl}{l_{\mathrm{Pl}}}
\newcommand{\RH}{R_{\mathrm{H}}}
\newcommand{\tH}{t_{\mathrm{H}}}
\newcommand{\ninf}{N_{\mathrm{e}}}
\newcommand{\ns}{n_{\mathrm{s}}}
\newcommand{\nt}{n_{\mathrm{T}}}
\newcommand{\rs}{r}
\newcommand{\alD}{\alpha_{\mathrm{D}}}
\newcommand{\beI}{\beta_{\mathrm{I}}}
\newcommand{\gaR}{\gamma_{\mathrm{R}}}
\newcommand{\Kcrit}{K_{\mathrm{crit}}}
\newcommand{\Rstruct}{R_{\mathrm{struct}}}
\newcommand{\Ntrials}{N_{\mathrm{trials}}}
\newcommand{\Nlife}{\langle N_{\mathrm{life}}\rangle}

% Operatoren
\DeclareMathOperator{\Tr}{Tr}
\DeclareMathOperator{\Var}{Var}
\DeclareMathOperator{\Cov}{Cov}
\DeclareMathOperator{\Div}{Div}
\DeclareMathOperator{\Grad}{Grad}

\hypersetup{
	colorlinks=true,
	linkcolor=blue,
	citecolor=blue,
	urlcolor=blue,
}

\begin{document}
	
	\begin{titlepage}
		\begin{center}
			\vspace*{1cm}
			
			\Huge\textbf{Anhang W: Passive Gravitationstomographie des Erdinneren unter Verwendung von Gezeitenwellen} \\
			\vspace{0.5cm}
			\Large\textbf{Ein koordinationsbasierter Ansatz innerhalb der Yakushev Unified Coordination Theory (YUCT)}
			
			\vspace{1.5cm}
			
			\Large\textbf{Alexey V. Yakushev\textsuperscript{1}}
			
			\vspace{0.3cm}
			\large
			\textsuperscript{1}Yakushev Research, \\
			YUCT-Theorie: \url{https://yuct.org/}\\
			YPSDC-Protokoll: \url{https://ypsdc.com/}\\
			
			\vspace{2cm}
			\textbf DOI: \url{https://doi.org/10.5281/zenodo.18444598}\\
			
			\vspace{1cm}
			
			\vspace{0cm}
			\large 
			Eine Kurzversion dieser Arbeit wurde bereits veröffentlicht und ist verfügbar unter\\ \url{https://doi.org/10.5281/zenodo.18362308}\\
			\vspace{1cm}
			März 2026 \\ \large V.2.0.
			
			\vspace{2cm}
			
			\begin{minipage}{0.9\textwidth}
				\begin{abstract}
					\noindent
					Dieser Anhang stellt eine neuartige Methode für die passive geophysikalische Exploration vor, die auf der Analyse von Resonanzen der Erdkruste unter dem Einfluss der Gezeitenkräfte von Mond und Sonne basiert. Im Rahmen der Yakushev Unified Coordination Theory (YUCT) zeigen wir, dass das Spektrum der Gezeitendeformationen Informationen über die unterirdische Verteilung der Koordinationseffizienz $\Keff$ enthält, die direkt mit Gesteinsart, Fluidgehalt (Öl, Gas, Wasser) und Erzkörpern zusammenhängt. Wir entwickeln ein mathematisches Modell, das gravimetrische und seismische Oberflächenmessungen mit der Tiefenverteilung von $\Keff$ verknüpft. Dabei nutzen wir das universelle Fehlerskalierungsgesetz $\eps = \kappac \alpha \Keff^{-\beta}$ ($\beta=0.67$, $\kappac=0.33$) sowie die Kopplungsgleichungen zwischen Gravitationsfeld und elastischen Deformationen aus den YUCT-Sektoren 0 (Gravitation) und 34 (Planetenwissenschaften). Wir zeigen, dass Gezeitenwellen eine kostengünstige, umweltfreundliche und global anwendbare Sondierungsmethode darstellen, deren Eindringtiefen mit aktiven Quellen nicht erreichbar sind. Wir geben einen Überblick über unabhängige Forschungsarbeiten, die Elemente dieses Ansatzes vorwegnehmen, einschließlich direkter Beobachtungen von Gezeitenresonanzen und Vorschlägen zur Kohlenwasserstoffexploration mittels Gravitationsgezeiten. Basierend auf einem koordinationsbasierten Mineral- und Fluid-Klassifikationssystem analog zum Periodensystem der Elemente wird eine Methodik zur Identifizierung von Kohlenwasserstoffen über charakteristische Resonanzfrequenzen und Gütefaktoren vorgeschlagen. Die wirtschaftliche Tragfähigkeit wird bewertet: Eine Verringerung des Bohrrisikos von Trockenlöchern um 20–30\% würde die Installationskosten um ein Mehrfaches amortisieren. Ferner erweitern wir die Methodik auf die Zustandsüberwachung von Bauwerken (Wolkenkratzer, Brücken, Staudämme), Erdbebenvorhersage und Tsunamifrüherkennung und zeigen, dass dieselben physikalischen Prinzipien eine revolutionäre neue Klasse passiver Überwachungstechnologien ermöglichen, die das Potenzial haben, Millionen von Leben und Billionen von Dollar zu retten. Die Ergebnisse eröffnen einen Weg zu einer grundlegend neuen Klasse von Explorations- und Überwachungstechnologien, die natürliche Gravitationsfelder nutzen.
				\end{abstract}
			\end{minipage}
			
			\vspace{1cm}
			
			\noindent
			\small\textbf{Schlüsselwörter:} YUCT, Gezeitenwellen, Koordinationseffizienz, Tiefentomographie, Kohlenwasserstoffexploration, Resonanzmethode, Geophysik, Planetare Innere, Zustandsüberwachung von Bauwerken, Erdbebenvorhersage, Tsunamifrüherkennung
			
			\vspace{0.5cm}
			
			\noindent
			\textcopyright 2026 Yakushev Research. Alle Rechte vorbehalten.
		\end{center}
	\end{titlepage}
	
	\tableofcontents
	\newpage
	
	\section{Einführung: Die Herausforderung der tiefen Geophysik}
	\label{sec:intro}
	
	\subsection{Die Grenzen der aktiven Exploration}
	
	Die Exploration von Kohlenwasserstoff- und Minerallagerstätten ist nach wie vor eine der kapitalintensivsten Aufgaben der angewandten Geowissenschaften. Konventionelle seismische Reflexionsmethoden erfordern leistungsstarke künstliche Quellen (Sprengstoffe, Vibroseis-LKW, Luftgewehre), was die Eindringtiefen auf 5–10 km begrenzt und erhebliche Umweltkosten verursacht. Gravimetrische und magnetische Vermessungen liefern nur tiefenintegrierte Informationen und können Fluidtypen nicht eindeutig identifizieren.
	
	\subsection{Gezeiten als natürliches Sondierungssignal}
	
	Die Erde wird kontinuierlich durch Gravitationskräfte von Mond und Sonne verformt, was zu Festkörpergezeiten mit Amplituden von bis zu 30–50 cm am Äquator führt. Diese Deformationen sind streng periodisch und weisen ein reiches harmonisches Spektrum auf, das halbtägige (12,42 h, 12,00 h), tägige (24–26 h) und langperiodische (14 d, 6 Monate, 18,6 a) Komponenten umfasst. Traditionell werden Gezeitensignale als Rauschen betrachtet, das bei hochpräzisen Messungen entfernt werden muss.
	
	Dieser Anhang schlägt einen Paradigmenwechsel vor: **Verwendung von Gezeitenwellen als natürliches, passives Sondierungssignal** zur Erforschung des Erdinneren. Im Rahmen der Yakushev Unified Coordination Theory (YUCT) zeigen wir, dass die Gezeitenantwort Informationen über die unterirdische Verteilung der Koordinationseffizienz $\Keff$ enthält, die als universeller Indikator für Gesteinseigenschaften, Fluidgehalt und Erzmineralisation dient.
	
	\subsection{Aufbau dieses Anhangs}
	
	Abschnitt \ref{sec:prior} gibt einen Überblick über unabhängige Forschungsarbeiten, die Elemente dieses Ansatzes vorwegnehmen und seine empirischen Grundlagen schaffen. Abschnitt \ref{sec:yuct-foundations} fasst die relevanten YUCT-Prinzipien zusammen. Abschnitt \ref{sec:model} entwickelt das mathematische Modell, das Gezeitendeformationen mit $\Keff$ verknüpft, und identifiziert explizit die beteiligten YUCT-Sektoren und Lagrange-Kopplungsterme. Abschnitt \ref{sec:identification} stellt ein koordinationsbasiertes Klassifikationssystem für geologische Materialien vor. Abschnitt \ref{sec:economics} bewertet die wirtschaftliche Tragfähigkeit. Abschnitt \ref{sec:roadmap} skizziert eine experimentelle Roadmap. Abschnitt \ref{sec:conclusion} diskutiert Implikationen für die Planetenforschung und zukünftige Missionen.
	
	\subsection{Von der Planetenforschung zur Zustandsüberwachung von Bauwerken}
	\label{sec:beyond-geophysics}
	
	Die in diesem Anhang entwickelte Methodik zur Untersuchung von Planeteninneren hat eine natürliche und leistungsstarke Erweiterung: die zerstörungsfreie Bewertung von menschengemachten Bauwerken. Jedes massive Objekt — ein Wolkenkratzer, eine Brücke, ein Staudamm, ein Kernreaktor-Sicherheitsbehälter — wird durch dieselben Gezeiten-Gravitationskräfte ständig verformt, die auch Himmelskörper formen. Diese Deformationen, obwohl winzig, tragen in ihrem Spektrum die Signatur des inneren Zustands des Objekts: seine Steifigkeitsverteilung, das Vorhandensein von Rissen, die Materialdegradation und die gesamte „Koordinationseffizienz“ \(K_{\mathrm{eff}}\) der Struktur.
	
	Diese Erkenntnis verwandelt die Gezeitentomographie von einem rein geophysikalischen Werkzeug in eine universelle Methode für die **passive Zustandsüberwachung von Bauwerken**. Im Gegensatz zu aktiven Methoden, die künstliche Anregung benötigen (Vibratoren, Schläge, Ultraschall), ist die Gezeitenüberwachung:
	\begin{itemize}
		\item \textbf{Passiv:} Nutzt natürliche, kostenlose und kontinuierliche Gravitationssignale.
		\item \textbf{Global:} Liefert ein integriertes Bild der gesamten Struktur, nicht nur lokaler Punkte.
		\item \textbf{Tief:} Durchdringt das gesamte Volumen und deckt versteckte innere Schäden auf.
		\item \textbf{Sicher:} Keine Strahlung, keine hochenergetischen Eingriffe, kein Kontakt erforderlich.
		\item \textbf{Kontinuierlich:} Ermöglicht Echtzeit- oder periodische Bewertung ohne Betriebsunterbrechung.
	\end{itemize}
	
	Die grundlegende Beziehung \(Q \propto K_{\mathrm{eff}}^{\beta}\) (Gl. \ref{eq:Q-keff}) verknüpft den beobachtbaren Gütefaktor einer Bauwerksresonanz direkt mit der Koordinationseffizienz des Materials. Mit zunehmendem Alter, Mikrorissen oder Schäden sinkt \(K_{\mathrm{eff}}\), und der Gütefaktor der Gezeitenresonanzen nimmt entsprechend ab. Die Überwachung dieser Veränderungen über die Zeit liefert eine quantitative, nicht-invasive Gesundheitsmetrik.
	
	\section{Unabhängige Vorarbeiten und empirische Grundlagen}
	\label{sec:prior}
	
	\subsection{Direkte Beobachtungen von Gezeitenresonanzen}
	
	Bemerkenswerterweise hat die Kernidee dieses Anhangs — dass Gezeitenwellen mit geologischen Strukturen resonant wechselwirken können — bereits direkte Beobachtungsbestätigung erhalten.
	
	\begin{quotation}
		\noindent
		\textit{„Resonanzen von langperiodischen (14–15 Tage) Gravitationsgezeiten und Deformationswellen in der Erdkruste, die durch die Verlagerung des Schwerpunkts des Erde-Mond-Systems auftraten, wurden in seismisch aktiven Regionen des Altai-Sajan, Kamtschatka und Sachalin während der Erdbebenüberwachung registriert... Der Artikel beweist, dass Resonanzen von Gravitationsgezeiten zur Erdbebenvorhersage sowie zur Frühwarnung vor geodynamischen Gefahren für Wasserbauwerke und \textbf{zur direkten Exploration von Öl- und Gasreservoiren} genutzt werden können.“} \cite{Sibgatulin2014}
	\end{quotation}
	
	Die Arbeit von Sibgatulin et al. \cite{Sibgatulin2014} liefert mehrere entscheidende empirische Grundlagen für die vorliegende Theorie:
	
	\begin{itemize}
		\item \textbf{Beobachtung von Resonanzen:} Langperiodische Gezeitenkomponenten (14–15 Tage) erzeugen messbare Resonanzantworten in der Kruste.
		\item \textbf{Auslösemechanismus:} Diese Resonanzen wirken in 80\% der Fälle als Auslöser für starke Erdbeben (M$\ge$5,0), was bestätigt, dass Gezeitenenergie effektiv an Krustenstrukturen koppelt.
		\item \textbf{Deformationswellen:} Die Resonanzen erzeugen langsame Deformationswellen (1–5 km/h), interpretiert als Solitonen, die sich durch die Kruste ausbreiten.
		\item \textbf{Explorationspotenzial:} Die Autoren schlagen explizit vor, diese Resonanzen für die direkte Kohlenwasserstoffexploration zu nutzen — eine vorausschauende Antizipation der hier entwickelten Methodik.
	\end{itemize}
	
	\subsection{Gezeitentomographie in der Planetenforschung}
	
	Der Begriff „Gezeitentomographie“ ist kürzlich in die planetologische Literatur eingegangen. Rovira-Navarro et al. \cite{RoviraNavarro2025} zeigen, dass:
	
	\begin{itemize}
		\item Das 3GM-Radioexperiment der JUICE-Mission gravitative Gezeitenvariationen des Ganymed mit ausreichender Genauigkeit messen wird, um laterale Variationen der inneren Struktur abzubilden.
		\item Schalendickenvariationen von 10–100\% erzeugbare Gezeitensignale erzeugen, wobei Kugelflächenfunktionen 2. und 3. Ordnung Äquator-Pol-Dickenunterschiede und hemisphärische Dichotomien einschränken.
		\item Die Autoren erklären ausdrücklich, dass „Gezeitentomographie“ — eine Technik, die zuvor auf der Erde angewendet wurde — auf Eismonde erweiterbar ist.
	\end{itemize}
	
	Diese Arbeit validiert das Grundprinzip, dass Gezeitenantworten auf die dreidimensionale innere Struktur empfindlich sind, und unterstützt den hier für die Erde entwickelten Ansatz.
	
	\subsection{Unterirdische Gezeitenmessungen}
	
	Rogister et al. \cite{Rogister2024} führten ein einzigartiges Experiment durch, bei dem sie Gezeiten-Gravitationssignale an der Oberfläche und in 520 m Tiefe im unterirdischen Labor Rustrel in Frankreich verglichen. Obwohl ihr Hauptaugenmerk auf der Gravimeterkalibrierung lag, zeigt ihre Arbeit, dass:
	
	\begin{itemize}
		\item Gezeiten-Gravitationsvariationen gleichzeitig in verschiedenen Tiefen mit supraleitenden Gravimetern gemessen werden können.
		\item Die theoretische Differenz zwischen Oberflächen- und Untergrundsignalen gering ist (\(8,5 \times 10^{-5}\) für halbtägige Wellen), aber prinzipiell nachweisbar.
		\item Aktuelle Kalibrierungsunsicherheiten die Auflösung dieser Differenz einschränken, aber verbesserte Instrumente tiefenabhängige Gezeitenmessungen ermöglichen könnten.
	\end{itemize}
	
	\subsection{Gezeitenantwort des Merkur und anderer Planeten}
	
	Briaud et al. \cite{Briaud2025} untersuchen die Gezeitenantwort des Merkur anhand von MESSENGER-Daten und bereiten sich auf BepiColombo-Beobachtungen vor. Sie betonen, dass:
	
	\begin{itemize}
		\item Gezeiten-Love-Zahlen empfindlich auf Inhomogenitäten reagieren, die aus räumlichen Variationen von Temperatur, Zusammensetzung und physikalischen Eigenschaften resultieren.
		\item Ähnlich wie bei Eismonden verändern lokale Temperatur- oder Mineralogieunterschiede die mechanischen Antworten, was zu einer nicht einheitlichen Gezeitendeformation führt.
		\item Die Messung langperiodischer Librationen, der Phasenverzögerung der Gezeiten und von Abweichungen vom Cassini-Zustand die innere Struktur einschränken kann.
	\end{itemize}
	
	Diese planetaren Anwendungen demonstrieren die Universalität der Gezeitentomographie als Werkzeug zur Untersuchung der inneren Struktur verschiedener Himmelskörper.
	
	\subsection{Synthese}
	
	Die oben zusammengefassten unabhängigen Forschungsarbeiten belegen:
	
	\begin{enumerate}
		\item \textbf{Gezeitenresonanzen sind beobachtbar} und korrelieren mit geologischen Strukturen \cite{Sibgatulin2014}.
		\item \textbf{Gezeitentomographie ist eine anerkannte Technik} in der Planetenforschung zur Kartierung dreidimensionaler innerer Strukturen \cite{RoviraNavarro2025}.
		\item \textbf{Unterirdische Gezeitenmessungen} sind mit aktueller Instrumentierung technisch machbar \cite{Rogister2024}.
		\item \textbf{Planetare Gezeitenantworten} sind empfindlich auf innere Inhomogenitäten \cite{Briaud2025}.
	\end{enumerate}
	
	Was bisher fehlte, ist ein einheitlicher theoretischer Rahmen, der diese Beobachtungen mit einem quantitativen Parameter ($\Keff$) verbindet, der für unterirdische Eigenschaften invertiert werden kann. Genau das liefert YUCT.
	
	\section{YUCT-Grundlagen für die Gezeitentomographie}
	\label{sec:yuct-foundations}
	
	\subsection{Koordinationseffizienz $\Keff$ und das universelle Fehlergesetz}
	
	Innerhalb von YUCT wird jedes physikalische System durch seine Koordinationseffizienz $\Keff$ charakterisiert, die den Grad der Kohärenz und Synchronisation seiner Komponenten quantifiziert. Das universelle Fehlerskalierungsgesetz (Anhang L) bezieht jeden relativen Fehler oder jede Fluktuation $\eps$ auf $\Keff$:
	
	\begin{equation}
		\eps = \kappac \, \alpha \, \Keff^{-\beta}, \qquad \beta = 0.67 \pm 0.03,\ \kappac = 0.33,
		\label{eq:universal-scaling}
	\end{equation}
	
	wobei $\alpha$ eine systemabhängige Konstante der Größenordnung $0.01$–$1$ ist. Dieses Gesetz wurde über 40 Größenordnungen hinweg verifiziert, von DNA-Replikationsfehlern bis zu Fluktuationen des kosmischen Mikrowellenhintergrunds.
	
	Im Kontext der Gezeitentomographie repräsentiert $\eps$ die relative Dämpfung oder Phasenverzögerung von Gezeitendeformationen, während $\Keff$ das geologische Medium (Gesteinstyp, Fluidgehalt, Porosität, Frakturierung) charakterisiert.
	
	\subsection{Für die Gezeitentomographie relevante YUCT-Sektorstruktur}
	
	Der YUCT-Lagrange-Formalismus (Anhang A) unterteilt die Realität in 120 wechselwirkende Sektoren. Für die Gezeitentomographie sind folgende Sektoren direkt relevant:
	
	\begin{itemize}
		\item \textbf{Sektor 0 (Gravitation):} Beschreibt die Metrik $g_{\mu\nu}$ und das Gravitationsfeld. Das von Mond und Sonne erzeugte Gezeitenpotential $W(\mathbf{r},t)$ tritt über diesen Sektor ein.
		\item \textbf{Sektor 34 (Planetenwissenschaften):} Beschreibt die innere Struktur von Planeten, einschließlich Dichte $\rho(\mathbf{r})$, elastischer Moduln $\lambda(\mathbf{r}), \mu(\mathbf{r})$ und Gezeitendeformationen.
		\item \textbf{Sektor 1 (Elektromagnetismus):} Relevant für magnetotellurische Korrelationen und elektromagnetische Vorläufer von Gezeitenresonanzen (wie von Sibgatulin et al. \cite{Sibgatulin2014} durch Messungen natürlicher pulsierender elektromagnetischer Felder beobachtet).
		\item \textbf{Sektor 62 (Ökologie) und 86 (Demographie):} Relevant für die Bewertung von Umwelt- und sozioökonomischen Auswirkungen von Explorationsaktivitäten (Abschnitt \ref{sec:economics}).
	\end{itemize}
	
	Die Kopplung zwischen Sektoren wird durch Koeffizienten $\kappa_{sr}$ vermittelt. Für die Gezeitentomographie ist der kritische Term:
	
	\begin{equation}
		\mathcal{L}_{\text{coupling}} = \kappa_{0,34} \, \mathrm{Tr}\left(\Psi_{0,34} \cdot O_0 \cdot O_{34}^\dagger\right),
		\label{eq:coupling}
	\end{equation}
	
	wobei $\Psi_{0,34}$ das Koordinationsfeld ist, das Gravitation und Planetenstruktur verbindet, $O_0$ den Gezeitenpotentialoperator und $O_{34}$ den Deformationstensoroperator darstellt.
	
	\subsection{Temperaturabhängigkeit und tiefe Struktur}
	
	Anhang P zeigt, dass $\Keff$ von der Temperatur abhängt, mit $\Keff(T) \propto T^{-\gamma}$ und $\beta\gamma = 0.20$, bestimmt aus geophysikalischen Daten (Erde-Mond-System, Weiße-Zwerge-Rotverschiebung). Diese Temperatursensitivität ist entscheidend für die Tiefentomographie, da geothermische Gradienten zusätzliche Randbedingungen für $K_{\mathrm{eff}}$-Profile liefern. Darüber hinaus bedeutet dies, dass die **thermische Modulation von Gezeitenresonanzen** genutzt werden kann, um Beiträge aus verschiedenen Tiefen zu trennen, da sich Oberflächentemperaturvariationen diffusiv ausbreiten und flachere Schichten schneller beeinflussen (siehe Abschn. \ref{sec:thermal-modulation}).
	
	\section{Mathematisches Modell der Gezeitenantwort}
	\label{sec:model}
	
	\subsection{Elastische Deformation unter Gezeitenpotential}
	
	Betrachten Sie ein elastisches Erdmodell mit Dichte $\rho(\mathbf{r})$, Lamé-Parametern $\lambda(\mathbf{r}), \mu(\mathbf{r})$, das dem Gezeitenpotential $W(\mathbf{r},t)$ externer Himmelskörper ausgesetzt ist. Die Bewegungsgleichung für kleine Verschiebungen $\mathbf{u}(\mathbf{r},t)$ lautet:
	
	\begin{equation}
		\rho \frac{\partial^2 \mathbf{u}}{\partial t^2} = \nabla \cdot \boldsymbol{\sigma} + \rho \nabla W,
		\label{eq:motion}
	\end{equation}
	
	wobei $\boldsymbol{\sigma}$ der Spannungstensor ist. Für ein isotropes elastisches Medium:
	
	\begin{equation}
		\sigma_{ij} = \lambda \delta_{ij} \nabla \cdot \mathbf{u} + \mu \left(\frac{\partial u_i}{\partial x_j} + \frac{\partial u_j}{\partial x_i}\right).
		\label{eq:stress}
	\end{equation}
	
	Das Gezeitenpotential wird nach Kugelflächenfunktionen entwickelt:
	
	\begin{equation}
		W(r,\theta,\phi,t) = \sum_{n=2}^{\infty} \sum_{m=-n}^{n} W_{nm}(r) Y_{nm}(\theta,\phi) e^{i\omega_{nm}t}.
		\label{eq:tidal-potential}
	\end{equation}
	
	Für die feste Erde dominieren die harmonischen Funktionen 2. Grades (halbtägige und tägige Wellen). Die Lösung wird als Entwicklung nach Kugelflächenfunktionen mit radialen Funktionen dargestellt, die das Love-Gleichungssystem erfüllen.
	
	\subsection{Einbeziehung von $\Keff$ in die elastischen Parameter}
	
	In YUCT werden elastische Eigenschaften durch $\Keff$ ausgedrückt. Nach der Methodik von Anhang C (Koordinationschemie) führen wir ein **koordinationsbasiertes Mineral-Klassifikationssystem** ein. Für ein gegebenes Gestein oder Fluid wird $\Keff$ durch empirisch kalibrierte Beziehungen mit dem Kompressionsmodul $K$, dem Schermodul $\mu$ und der Dichte $\rho$ verknüpft:
	
	\begin{align}
		\lambda(\mathbf{r}) &= \lambda_0 \cdot f_\lambda(\Keff(\mathbf{r})), \\
		\mu(\mathbf{r}) &= \mu_0 \cdot f_\mu(\Keff(\mathbf{r})), \\
		\rho(\mathbf{r}) &= \rho_0 \left(\frac{\Keff(\mathbf{r})}{\Keffzero}\right)^{\gamma_\rho}, \quad \gamma_\rho \approx 0.2\ \text{(Bereich 0,1–0,3)},
		\label{eq:keff-properties}
	\end{align}
	
	Die Funktionen $f_\lambda, f_\mu$ sind monoton wachsend, was widerspiegelt, dass eine höhere Koordination (kohärentere atomare/molekulare Struktur) zu einer größeren Steifigkeit führt. Vorläufige Schätzungen deuten auf $f_\lambda(\Keff) \propto \Keff^{\xi}$ mit $\xi \approx 0.3$–$0.5$ hin, basierend auf der Skalierung elastischer Moduln mit Dichte und Koordinationszahl.
	
	\subsection{Resonanzantwort und Gütefaktor}
	
	Wenn die Frequenz $\omega$ einer Gezeitenharmonischen sich der Eigenfrequenz $\omega_0$ einer unterirdischen Struktur (z.B. eines fluidgefüllten Reservoirs) nähert, tritt eine resonante Verstärkung auf. Der Gütefaktor $Q$ der Resonanz ist mit $\Keff$ über das universelle Fehlergesetz verbunden. Die Energiedissipation pro Zyklus ist proportional zum relativen Fehler $\eps$, und da $1/Q = \Delta E / (2\pi E) \approx \eps$, gilt:
	
	\begin{equation}
		\frac{1}{Q} \propto \eps = \kappac \alpha \Keff^{-\beta} \quad \Longrightarrow \quad Q = Q_0 \, \Keff^{\beta},
		\label{eq:Q-keff}
	\end{equation}
	
	wobei $Q_0$ ein Referenzwert ist. Somit:
	\begin{itemize}
		\item Medien mit hohem $\Keff$ (kompetente Gesteine, dichte Fluide) erzeugen scharfe, hoch-$Q$-Resonanzen.
		\item Medien mit niedrigem $\Keff$ (frakturierte Zonen, Gas) erzeugen breite, niedrig-$Q$-Merkmale.
	\end{itemize}
	
	\subsection{Vorwärtsmodell und inverses Problem}
	
	Oberflächenmessungen von Gravitationsvariationen und Verschiebungen liefern Zeitreihen $d(t)$. Nach Extraktion der Gezeitenharmonischen (durch Faltung mit bekannten Ephemeriden) und Subtraktion der Antwort eines radial symmetrischen Referenzmodells enthält die Residue $\delta d(t)$ Informationen über laterale Inhomogenitäten.
	
	Die Fourier-Transformation ergibt das Spektrum $\delta D(\omega)$, das Peaks bei Frequenzen $\omega_i$ enthält, die Resonanzen unterirdischer Strukturen entsprechen. Aus jedem Peak extrahieren wir:
	
	\begin{itemize}
		\item $\omega_i$ — Resonanzfrequenz, empfindlich auf Tiefe, Größe und elastische Eigenschaften.
		\item $Q_i = \omega_i / \Delta \omega_i$ — Gütefaktor, verknüpft mit $\Keff$ über Gl. (\ref{eq:Q-keff}).
		\item $A_i$ — Amplitude, proportional zum $\Keff$-Kontrast zwischen Struktur und Hintergrund.
	\end{itemize}
	
	Das inverse Problem — Wiederherstellung des dreidimensionalen $\Keff(\mathbf{r})$ — wird durch tomographische Inversion gelöst, die die Diskrepanz zwischen beobachteten und synthetischen Spektren minimiert. Die Regularisierung nutzt Vorabinformationen aus dem koordinationsbasierten Mineral-Klassifikationssystem.
	
	\subsection{Thermische Modulation von Gezeitenresonanzen}
	\label{sec:thermal-modulation}
	
	Temperaturschwankungen beeinflussen $\Keff$ über die Beziehung $\Keff(T) \propto T^{-\gamma}$ (Anhang P). Tägliche und jahreszeitliche Temperaturänderungen breiten sich als thermische Wellen in den Untergrund aus und modulieren das lokale $\Keff$ und damit die Resonanzfrequenzen und Gütefaktoren flacher Strukturen. Dies bietet eine zusätzliche Dimension für die Tiefendiskriminierung: Tiefere Strukturen reagieren mit einer Phasenverzögerung und reduzierter Amplitude auf die Oberflächentemperaturerregung.
	
	Wenn das Temperaturfeld $T(\mathbf{r},t)$ aus meteorologischen Daten bekannt oder modelliert ist, kann die zeitabhängige Störung $\delta \Keff(\mathbf{r},t) \approx -\gamma \Keff(\mathbf{r}) \delta T(\mathbf{r},t)/T$ berechnet werden. Diese moduliert die Gezeitenantwort und erzeugt Seitenbänder oder Amplitudenmodulation bei thermischen Frequenzen. Die Analyse dieser Modulationen ermöglicht die Trennung flacher und tiefer Beiträge und verbessert die vertikale Auflösung.
	
	\subsection{Erweiterung auf technische Bauwerke: Wolkenkratzer, Brücken und Dämme}
	\label{sec:structures}
	
	Das in den Abschnitten \ref{sec:model} entwickelte mathematische Modell ist direkt auf große menschengemachte Strukturen anwendbar, mit geeigneten Modifikationen der Randbedingungen und Geometrie. Für eine Struktur mit Masse \(M\), Höhe \(H\) und charakteristischer Steifigkeit \(k\) skaliert die Grundfrequenz der Biegeschwingung wie folgt:
	
	\begin{equation}
		f_0 \approx \frac{1}{2\pi} \sqrt{\frac{k}{M_{\mathrm{eff}}}} \sim \frac{1}{H} \sqrt{\frac{E}{\rho}},
		\label{eq:structural-frequency}
	\end{equation}
	
	wobei \(E\) der Elastizitätsmodul, \(\rho\) die Dichte und \(M_{\mathrm{eff}}\) die effektive modale Masse ist. Für einen typischen Wolkenkratzer (\(H \sim 300\) m, \(E/\rho \sim 10^7\) m\(^2\)/s\(^2\)) gilt \(f_0 \sim 0.1\) Hz — nach nichtlinearer Mischung oder parametrischer Erregung gut innerhalb des Bereichs der Gezeitenharmonischen.
	
	Der Gütefaktor \(Q\) dieser Moden gehorcht derselben universellen Beziehung:
	
	\begin{equation}
		Q = Q_0 \left( \frac{K_{\mathrm{eff}}}{K_{\mathrm{eff},0}} \right)^{\beta},
		\label{eq:structural-Q}
	\end{equation}
	
	wobei \(K_{\mathrm{eff}}\) nun die **strukturelle Koordinationseffizienz** bezeichnet — ein dimensionsloses Maß dafür, wie gut alle Komponenten der Struktur zusammenarbeiten. Neu errichtete Gebäude haben ein hohes \(K_{\mathrm{eff}}\); beschädigte oder degradierte Strukturen weisen niedrigere Werte auf.
	
	\subsubsection{Was die Gezeitenüberwachung offenbart}
	
	Durch die Platzierung hochempfindlicher Beschleunigungsmesser oder Gravimeter an einer Struktur (z.B. auf dem Dach, in halber Höhe und am Fundament) und Aufzeichnung über mehrere Wochen erhalten wir:
	
	\begin{itemize}
		\item \textbf{Modale Frequenzen:} Verschiebungen weisen auf Steifigkeitsänderungen hin (z.B. durch Rissbildung oder Materialdegradation).
		\item \textbf{Gütefaktoren:} Abnahmen signalisieren akkumulierte Schäden, Mikrorisse oder gelockerte Verbindungen.
		\item \textbf{Anisotropie:} Unterschiede in der Antwort entlang verschiedener Achsen offenbaren gerichtete Schadensmuster.
		\item \textbf{Höhere Moden:} Liefern tiefenaufgelöste Informationen über die Schadenslokalisierung.
	\end{itemize}
	
	\subsubsection{Das Konzept des „gravitativen Fingerabdrucks“}
	
	Jede große Struktur hat einen einzigartigen „gravitativen Fingerabdruck“ — ihr Spektrum aus Resonanzfrequenzen und Gütefaktoren, das durch Gezeitenkraft angeregt wird. Dieser Fingerabdruck kann:
	\begin{enumerate}
		\item \textbf{Bei Inbetriebnahme aufgezeichnet} werden, um eine Basislinie zu erstellen.
		\item \textbf{Regelmäßig überwacht} werden, um die Degradation zu verfolgen.
		\item \textbf{Nach extremen Ereignissen} (Erdbeben, Hurrikane, Brände) erneut gemessen werden, um Schäden ohne Sichtprüfung zu bewerten.
	\end{enumerate}
	
	Die Empfindlichkeit ist bemerkenswert: Eine 1\%-Änderung der Steifigkeit (z.B. durch einen großen Riss) verschiebt die Frequenzen um \(\sim 0.5\%\), was mit modernen Instrumenten leicht nachweisbar ist.
	
	\subsection{Erdbebenvorhersage: Überwachung kritischer Verwerfungen mit Gezeitentomographie}
	\label{sec:earthquakes}
	
	Dieselben Gezeitenkräfte, die die Tomographie von Planeteninneren und technischen Bauwerken ermöglichen, belasten auch kontinuierlich die Erdkruste. Seit Jahrzehnten beobachten Seismologen statistische Korrelationen zwischen Gezeitenphase und Erdbebenauftreten, insbesondere für Überschiebungs- und Abschiebungsereignisse \cite{Tanaka2012, Tolstoy2013, Sibgatulin2014, RAN2025}. Diese Korrelationen, obwohl signifikant (bis zu 10\% Wahrscheinlichkeitserhöhung bei bestimmten Mondabständen), sind beschreibend geblieben — ein Phänomen ohne Vorhersagetheorie.
	
	YUCT liefert den fehlenden Rahmen. Eine geologische Verwerfung, die sich dem Versagen nähert, ist ein System, dessen Koordinationseffizienz \(K_{\mathrm{eff}}\) mit zunehmender Mikrorissbildung und Spannungskonzentration abnimmt. Das universelle Fehlerskalierungsgesetz \(\varepsilon = \kappa_c \alpha K_{\mathrm{eff}}^{-\beta}\) (Gl. \ref{eq:universal-scaling}) bestimmt, wie das System auf Störungen — hier Gezeitenkräfte — reagiert.
	
	\subsubsection{Mechanismus}
	
	\begin{enumerate}
		\item \textbf{Spannungsakkumulation im Hintergrund:} Plattenbewegungen laden die Verwerfung über Jahre bis Jahrhunderte auf.
		\item \textbf{Kritischer Zustand:} Wenn die Spannung sich der Versagensschwelle nähert, sinkt \(K_{\mathrm{eff}}\), wodurch die Verwerfung zunehmend empfindlich auf kleine Störungen wird.
		\item \textbf{Gezeitenauslösung:} Das periodische Gezeitenpotential \(W(\mathbf{r},t)\) von Mond und Sonne addiert eine kleine oszillatorische Spannung. Wenn diese Spannung mit der Verwerfungsebene übereinstimmt und eine kritische Amplitude erreicht, kann sie die Verwerfung lösen.
		\item \textbf{Versagen:} Das Erdbeben tritt ein.
	\end{enumerate}
	
	Die Auslösewahrscheinlichkeit skaliert wie:
	\[
	P_{\text{trigger}}(t) \propto \left( \frac{W(t)}{W_0} \right)^{\beta}, \quad \beta = 0.67,
	\]
	eine direkte Folge des universellen Gesetzes. Diese nichtlineare Skalierung erklärt, warum statistische Korrelationen signifikant, aber nicht deterministisch sind — die Antwort der Verwerfung wird durch ihr sinkendes \(K_{\mathrm{eff}}\) verstärkt.
	
	\subsubsection{Was die Gezeitentomographie hinzufügt}
	
	Ein Netzwerk hochempfindlicher Gravimeter, das um eine bekannte Verwerfungszone herum eingesetzt wird, ermöglicht:
	
	\begin{itemize}
		\item \textbf{Kontinuierliche Überwachung von \(K_{\mathrm{eff}}\):} Die Antwort der Verwerfung auf jede Gezeitenharmonische (z.B. M2, O1, Mf) liefert eine Echtzeitschätzung ihrer Koordinationseffizienz. Ein anhaltender Abfall signalisiert eine zunehmende Gefahr.
		\item \textbf{Räumliche Kartierung:} Mehrere Stationen können die Region des abnehmenden \(K_{\mathrm{eff}}\) lokalisieren und identifizieren, welches Segment der Verwerfung sich der Kritikalität nähert.
		\item \textbf{Auslösevorhersage:} Durch Kombination des Echtzeit-\(K_{\mathrm{eff}}\) mit präzisen Ephemeriden können Zeitfenster erhöhten Risikos vorhergesagt werden — Tage, an denen die Gezeitenspannung mit der Verwerfung übereinstimmt und die Auslöseschwelle überschreitet.
	\end{itemize}
	
	\subsubsection{Empirische Unterstützung}
	
	Mehrere unabhängige Studien haben die zugrunde liegenden Korrelationen bereits nachgewiesen:
	
	\begin{itemize}
		\item Tanaka (2012) zeigte, dass gezeitenausgelöste Mikrobeben im Jahrzehnt vor dem Tohoku-oki-Erdbeben 2011 zunahmen und danach verschwanden \cite{Tanaka2012}.
		\item Tolstoy (2013) dokumentierte, dass der Meeresboden mit den Gezeiten „atmet“ und die Mikroseismizität auf mittelozeanischen Rücken mit der Gezeitenphase korreliert \cite{Tolstoy2013}.
		\item Sibgatulin et al. (2014) fanden, dass 80\% der starken Erdbeben mit 14-tägigen Gezeitenresonanzen korrelieren \cite{Sibgatulin2014}.
		\item Jüngste Studien der Russischen Akademie der Wissenschaften (2025) zeigen eine Wahrscheinlichkeitserhöhung von 9–10\% für Erdbeben bei bestimmten Erde-Mond-Abständen \cite{RAN2025}.
		\item Deformationsvorläufer: Stufenartige Volumendehnungsänderungen Tage bis Wochen vor moderaten Erdbeben wurden Gezeitenkräften zugeschrieben \cite{DeformationPrecursors2024}.
	\end{itemize}
	
	Was bisher fehlte, war ein einheitlicher theoretischer Rahmen, der diese Beobachtungen mit einem messbaren physikalischen Parameter verbindet. YUCT liefert diese Verbindung: \(K_{\mathrm{eff}}\) der Verwerfungszone.
	
	\subsubsection{Praktische Umsetzung}
	
	\begin{itemize}
		\item \textbf{Instrumentierung:} 10–20 supraleitende Gravimeter, die um eine gut untersuchte Verwerfung herum eingesetzt werden (z.B. San-Andreas-Verwerfung, Nordanatolische Verwerfung, Japan-Graben).
		\item \textbf{Dauer:} Kontinuierlicher Betrieb über 5–10 Jahre, um mehrere Erdbebenzyklen zu erfassen.
		\item \textbf{Erwartetes Ergebnis:} Retrospektive Analysen sollten vor großen Ereignissen abfallende \(K_{\mathrm{eff}}\)-Vorläufer zeigen. Prospektive Überwachung könnte dann probabilistische Warnungen ausgeben.
	\end{itemize}
	
	Die Kosten eines solchen Netzwerks (1–2 Millionen \$) sind im Vergleich zu den wirtschaftlichen und menschlichen Kosten eines einzigen großen Erdbebens vernachlässigbar. Eine 24-Stunden-Warnung für ein Ereignis der Stärke 7+ könnte Zehntausende von Leben und Milliarden von Dollar retten.
	
	\subsubsection{Quantitative Vorhersage: Ein \(K_{\mathrm{eff}}\)-Abfall als Vorläufer}
	\label{subsubsec:precursor}
	
	Ein Schlüsselelement der Erdbebenvorhersage ist die Fähigkeit, Veränderungen von \(K_{\mathrm{eff}}\) im Laufe der Zeit entlang einer Verwerfung zu beobachten. Nach dem universellen Fehlerskalierungsgesetz ist die relative Dehnung (oder Mikrorissdichte) \(\varepsilon\) mit \(K_{\mathrm{eff}}\) verknüpft: \(\varepsilon = \kappa_c \alpha K_{\mathrm{eff}}^{-\beta}\). Wenn sich tektonische Spannungen aufbauen, nimmt die Koordinationseffizienz der Verwerfungszone ab, was zu einer Zunahme der Dehnung und folglich zu einer Zunahme der Amplitude und einer Änderung der Phase der Gezeitenantwort führt.
	
	Betrachten Sie ein Modellszenario. Nehmen wir an, dass zu einem Zeitpunkt \(t_0\) (Hintergrundzustand) die Koordinationseffizienz der Verwerfung \(K_{\mathrm{eff}}^{(0)}\) ist. Wenn sie sich einem kritischen Zustand nähert, fällt sie um \(\Delta K\). Die relative Änderung der Amplitude der Gezeitenantwort kann aus Gleichung (25) abgeschätzt werden:
	
	\[
	\frac{\Delta A}{A} \approx \frac{\Delta Q}{Q} \approx \beta \frac{\Delta K}{K}.
	\]
	
	Für \(\beta = 0.67\) und \(\Delta K/K = 20\%\) erhalten wir \(\Delta A/A \approx 13\%\). Eine solche Amplitudenänderung ist mit modernen Gravimetern durch monatliche Mittelung leicht nachweisbar.
	
	Darüber hinaus kann der Abfall von \(K_{\mathrm{eff}}\) selbst mit der Auslösewahrscheinlichkeit verknüpft werden. Aus Gleichung (\ref{eq:trigger}) folgt, dass die Wahrscheinlichkeit, dass ein Erdbeben durch eine Gezeitenwelle ausgelöst wird, proportional zu \((W(t)/W_0)^{\beta}\) ist. Mit abnehmendem \(K_{\mathrm{eff}}\) sinkt folglich die Schwellenamplitude \(W_0\) (d.h. die minimale Gezeitenbelastung, die ein Ereignis auslösen kann). In der Praxis bedeutet dies: Wenn ein Gravimeternetzwerk einen stetigen Abfall von \(K_{\mathrm{eff}}\) in einer bestimmten Verwerfungszone um 20\% oder mehr über mehrere Wochen aufzeichnet, dann:
	
	\begin{enumerate}
		\item Das Risiko eines Erdbebens in den kommenden Tagen steigt.
		\item Die Zeiten, in denen Gezeitenwellen ihre maximale Amplitude erreichen (z.B. M2, Mf), werden zu Fenstern erhöhter Wahrscheinlichkeit.
		\item Wenn \(K_{\mathrm{eff}}\) unter einen kritischen Schwellenwert \(K_{\text{crit}}\) fällt, kann eine Warnung ausgegeben werden.
	\end{enumerate}
	
	Somit ermöglicht die kontinuierliche Überwachung von \(K_{\mathrm{eff}}\) mittels Daten eines Gravimeternetzwerks nicht nur die Diagnose des aktuellen Zustands einer Verwerfung, sondern auch die kurzzeitige (Tage bis Wochen) Vorhersage der Wahrscheinlichkeit starker Ereignisse. Dies ist eine direkte praktische Anwendung von YUCT in der Geodynamik.
	
	\subsection{Tsunamivorhersage: Die ultimative Anwendung}
	\label{sec:tsunami}
	
	Unterseebeben, die Tsunamis verursachen, sind die gefährlichsten und am wenigsten vorhersagbaren aller Naturkatastrophen. Während 80\% der Tsunamis durch koseismische Deformation des Meeresbodens verursacht werden \cite{NOAA2024, USGS2024}, aktivieren sich die aktuellen Warnsysteme erst \textit{nach} dem Erdbeben. Die kostbaren Minuten zwischen dem Bruch und dem Eintreffen der Welle werden für Berechnungen aufgewendet, nicht für die Evakuierung.
	
	Die Gezeitentomographie bietet einen Paradigmenwechsel: **Vorhersage des Bruchs, bevor er geschieht**.
	
	\subsubsection{Mechanismus}
	
	Eine Subduktionszone, die sich dem Versagen nähert, verhält sich wie ein System mit abnehmender Koordinationseffizienz \(K_{\mathrm{eff}}\). Mit zunehmender Spannungsakkumulation nehmen Mikrorisse zu, und die Antwort der Verwerfung auf die Gezeitenkräfte wird zunehmend nichtlinear. Das universelle Skalierungsgesetz bestimmt diese Antwort:
	\[
	\varepsilon = \kappa_c \alpha K_{\mathrm{eff}}^{-\beta}, \quad \beta = 0.67.
	\]
	
	Die Überwachung eines Netzwerks von Meeresboden-Gravimetern um eine Subduktionszone (z.B. Japan-Graben, Cascadia, Sumatra) liefert:
	\begin{itemize}
		\item \textbf{Echtzeit-\(K_{\mathrm{eff}}\)-Schätzung:} Die Amplitude und Phase der Gezeitenharmonischen (M2, O1, Mf) werden kontinuierlich gemessen und auf die Verwerfungskoordination zurückgerechnet.
		\item \textbf{Vorläufererkennung:} Ein anhaltender Abfall von \(K_{\mathrm{eff}}\) über Wochen bis Monate signalisiert, dass die Verwerfung in einen kritischen Zustand eintritt.
		\item \textbf{Auslösevorhersage:} Durch Kombination des aktuellen \(K_{\mathrm{eff}}\) mit präzisen Ephemeriden können Zeitfenster erhöhter Wahrscheinlichkeit vorhergesagt werden:
		\[
		P_{\text{trigger}}(t) \propto \left( \frac{W(t)}{W_0} \right)^{0.67},
		\]
		wobei \(W(t)\) das Gezeitenpotential zum Zeitpunkt \(t\) ist.
	\end{itemize}
	
	\subsubsection{Empirische Unterstützung}
	
	Mehrere unabhängige Studien haben die zugrunde liegenden Korrelationen bereits nachgewiesen:
	\begin{itemize}
		\item Tanaka (2012) zeigte, dass gezeitenausgelöste Mikrobeben im Jahrzehnt vor dem Tohoku-oki-Erdbeben 2011 zunahmen und danach verschwanden \cite{Tanaka2012}.
		\item Tolstoy (2013) dokumentierte, dass der Meeresboden mit den Gezeiten „atmet“ und die Mikroseismizität auf mittelozeanischen Rücken mit der Gezeitenphase korreliert \cite{Tolstoy2013}.
		\item Sibgatulin et al. (2014) fanden, dass 80\% der starken Erdbeben mit 14-tägigen Gezeitenresonanzen korrelieren \cite{Sibgatulin2014}.
		\item Jüngste Studien der Russischen Akademie der Wissenschaften (2025) zeigen eine Wahrscheinlichkeitserhöhung von 9–10\% für Erdbeben bei bestimmten Erde-Mond-Abständen \cite{RAN2025}.
	\end{itemize}
	
	Was bisher fehlte, war ein einheitlicher theoretischer Rahmen, der diese Beobachtungen mit einem messbaren physikalischen Parameter verbindet. YUCT liefert diesen Rahmen: \(K_{\mathrm{eff}}\) der Verwerfungszone.
	
	\subsubsection{Praktische Umsetzung}
	
	\begin{itemize}
		\item \textbf{Instrumentierung:} 20–30 Meeresboden-Gravimeter (Kosten jeweils ca. 50.000 \$), die in einem 200×200 km großen Gitter über einer bekannten Subduktionszone eingesetzt werden.
		\item \textbf{Dauer:} Kontinuierlicher Betrieb über 5–10 Jahre, um den gesamten Erdbebenzyklus zu erfassen.
		\item \textbf{Datenfusion:} Kombinieren Sie gravimetrische Daten mit bestehenden DART-Bojen-Netzwerken \cite{NOAA2024} und seismischer Überwachung \cite{USGS2024} für eine Multi-Parameter-Validierung.
		\item \textbf{Warnprotokoll:} Wenn \(K_{\mathrm{eff}}\) unter einen kalibrierten Schwellenwert fällt und mit einem Flut-Hochwasserfenster zusammenfällt, werden automatisch Warnungen an die Katastrophenschutzbehörden ausgegeben.
	\end{itemize}
	
	Die Kosten eines solchen Netzwerks (1–2 Millionen \$) sind im Vergleich zu den wirtschaftlichen und menschlichen Kosten eines einzigen großen Tsunamis vernachlässigbar. Der Indische-Ozean-Tsunami 2004 tötete 250.000 Menschen und verursachte Milliardenschäden \cite{NOAA2024, USGS2024}. Eine 24-Stunden-Warnung vor einem ähnlichen Ereignis wäre die größte humanitäre Errungenschaft des 21. Jahrhunderts.
	
	\subsubsection{Fazit}
	
	Die Tsunamivorhersage ist kein Traum — sie ist eine direkte Konsequenz von YUCT. Indem wir messen, wie der Meeresboden mit den Gezeiten „atmet“ \cite{Tolstoy2013}, können wir endlich sehen, wann die Erde im Begriff ist, Zerstörung auszuatmen. Dies ist keine Spekulation; es ist Physik.
	
	\section{Identifizierung von Kohlenwasserstoffen und Mineralien mittels \(\Keff\)-Signaturen}
	\label{sec:identification}
	
	\subsection{Koordinationsbasierte Klassifikation geologischer Materialien}
	
	Analog zum Periodensystem der Elemente (Anhang C) können geologische Materialien anhand ihrer charakteristischen \(\Keff\)-Werte klassifiziert werden. Tabelle \ref{tab:keff-geology} gibt vorläufige Schätzungen basierend auf koordinationschemischen Prinzipien und Skalierungsbeziehungen. Die Quellen für diese Schätzungen umfassen:
	\begin{itemize}
		\item Elastische Moduln und Dichten aus Standardreferenzwerken (z.B. Carmichael, "Practical Handbook of Physical Properties of Rocks and Minerals", CRC Press).
		\item Empirische Korrelationen zwischen \(\Keff\) und chemischer Zusammensetzung aus Anhang C.
		\item Inversion von seismischen und Bohrlochdaten aus gut untersuchten Ölfeldern (vorläufig).
	\end{itemize}
	
	\begin{table}[htbp]
		\centering
		\caption{Koordinationseffizienz \(\Keff\) für repräsentative geologische Materialien (vorläufige Schätzungen)}
		\label{tab:keff-geology}
		\small
		\begin{tabular}{lcc}
			\toprule
			\textbf{Material} & \textbf{\(\Keff\) (geschätzt)} & \textbf{Bemerkungen} \\
			\midrule
			\multicolumn{3}{l}{\textit{Fluide}} \\
			Erdgas (Methan) & $6.1 \pm 0.3$ & Hohe Beweglichkeit, niedrige Dichte \\
			Leichtes Rohöl & $5.2 \pm 0.2$ & Mittlere Koordination \\
			Schweres Rohöl & $4.8 \pm 0.2$ & Höhere Viskosität, niedrigeres \(\Keff\) \\
			Wasser (Formationswasser) & $4.5 \pm 0.2$ & Starke Wasserstoffbrückenbindungen \\
			\hline
			\multicolumn{3}{l}{\textit{Gesteine}} \\
			Unverfestigter Sand & $3.2 \pm 0.3$ & Niedrige Koordination, hohe Porosität \\
			Sandstein (typisch) & $3.8 \pm 0.3$ & Mittel \\
			Sandstein (dicht) & $4.2 \pm 0.3$ & Reduzierte Porosität \\
			Kalkstein & $4.5 \pm 0.3$ & Kristalline Struktur \\
			Granit & $5.0 \pm 0.4$ & Hohe Kohärenz \\
			Basalt & $5.3 \pm 0.4$ & Dicht, feinkörnig \\
			\hline
			\multicolumn{3}{l}{\textit{Erzmineralien}} \\
			Gold & $8.2 \pm 0.5$ & Hohe Dichte, metallische Bindung \\
			Kupfer & $7.5 \pm 0.5$ & Metallisch \\
			Eisenerz (Hämatit) & $6.8 \pm 0.4$ & Oxidisch \\
			\bottomrule
		\end{tabular}
	\end{table}
	
	\subsection{Resonanzsignaturen von Kohlenwasserstoffreservoiren}
	
	Ein Kohlenwasserstoffreservoir besteht typischerweise aus porösem Speichergestein (Sandstein, Karbonat), dessen Porenraum mit Öl oder Gas gefüllt ist, und wird durch eine undurchlässige Deckschicht (Tonstein, Evaporit) versiegelt. Diese Konfiguration erzeugt eine resonante Struktur mit einer charakteristischen Frequenz, die bestimmt wird durch:
	
	\begin{equation}
		\omega_0 \approx \frac{\pi v_s}{2h} \sqrt{\frac{\rho_f}{\rho_r} \cdot \frac{K_r}{K_f}},
		\label{eq:resonant-frequency}
	\end{equation}
	
	wobei $v_s$ die Scherwellengeschwindigkeit, $h$ die Reservoir-Mächtigkeit, $\rho_f, \rho_r$ die Fluiddichten und die Gesteinsdichte sowie $K_f, K_r$ die Kompressionsmoduln sind. Dieser Ausdruck leitet sich von der Grundmode einer in eine elastische Matrix eingebetteten Fluid-Schicht ab (die „Quetsch“-Resonanz). Für typische Reservoirparameter ($v_s \approx 2500$ m/s, $h = 50$ m, $\rho_f/\rho_r \approx 0.8$, $K_r/K_f \approx 10$) ist $\omega_0/(2\pi) \approx 0.5$ Hz — nach Skalierung noch innerhalb des Gezeitenbandes? Tatsächlich sind Gezeitenfrequenzen viel niedriger ($\sim 10^{-5}$ Hz). Allerdings können Resonanzen bei viel niedrigeren Frequenzen auftreten, wenn die effektive Nachgiebigkeit hoch ist (z.B. großes Fluidvolumen). Alternativ kann die nichtlineare Mischung von Gezeitenfrequenzen höherfrequente Harmonische erzeugen. Dies bleibt ein Bereich für weitere Forschung.
	
	Der Gütefaktor $Q$ ist durch Gl. (\ref{eq:Q-keff}) gegeben, wobei $\Keff$ die effektive Koordination des Reservoirs repräsentiert. Unter Verwendung von Tabelle \ref{tab:keff-geology}:
	
	\begin{itemize}
		\item Gasreservoir: $\Keff \approx 6.1 \Rightarrow Q \approx Q_0 \times 6.1^{0.67} \approx 3.5\,Q_0$ (hoher $Q$, scharfe Resonanz)
		\item Ölreservoir: $\Keff \approx 5.2 \Rightarrow Q \approx Q_0 \times 5.2^{0.67} \approx 3.0\,Q_0$ (mittel)
		\item Wassergesättigt: $\Keff \approx 4.5 \Rightarrow Q \approx Q_0 \times 4.5^{0.67} \approx 2.7\,Q_0$ (niedrigerer $Q$)
	\end{itemize}
	
	Somit kann die Spektralanalyse von Gezeitenresonanzen zwischen gas-, öl- und wasserführenden Strukturen unterscheiden, ohne zu bohren.
	
	% ========= KORRIGIERTER ABSCHNITT =========
	\subsubsection{Erwartete Signalamplitude und Messbarkeit durch moderne Gravimeter}
	\label{subsubsec:signal_amplitude}
	
	Um die praktische Machbarkeit der Methode zu bewerten, muss gezeigt werden, dass die Amplitude der Gezeitenresonanz eines typischen Öl- und Gasreservoirs die Empfindlichkeitsschwelle moderner Gravimeter (\( \sim 10^{-11}\ \text{m/s}^2\)) überschreitet. Betrachten Sie ein Reservoir mit einer Mächtigkeit \(h = 50\) m, gasgesättigt (\(\Keff \approx 6.1\)), das in einer Tiefe \(H = 2000\) m in einem homogenen Wirtsgestein mit \(\Keff \approx 4.0\) liegt. Der Kontrast in der Koordinationseffizienz \(\Delta \Keff / \Keff \approx 0.35\) erzeugt einen entsprechenden Kontrast in Dichte und Steifigkeit, der über die Zustandsgleichung abgeschätzt werden kann.
	
	Die M2-Gezeitenwelle (Periode 12,42 h) erzeugt an der Erdoberfläche Gravitationsvariationen mit einer Amplitude von etwa 100–150 nm/s\(^2\) (d.h. \(\sim 10^{-8} g\) je nach Breitengrad). Die resonante Verstärkung im Reservoir führt zu einer zusätzlichen lokalen Anomalie \(\delta g_{\text{res}}\), die wie folgt geschätzt werden kann:
	
	\begin{equation}
		\delta g_{\text{res}} \approx \delta g_{\text{tide}} \cdot \frac{Q}{Q_0} \cdot \frac{\Delta \Keff}{\Keff} \cdot \left(\frac{h}{H}\right)^2,
		\label{eq:signal_estimate}
	\end{equation}
	
	wobei \(\delta g_{\text{tide}} \approx 100\) nm/s\(^2\) das Hintergrund-Gezeitensignal ist, \(Q\) der Gütefaktor der Resonanz (aus Gleichung (25)), \(Q_0 \approx 100\) der Gütefaktor der Hintergrundoszillationen und \(h/H\) ein geometrischer Faktor ist, der die Tiefe berücksichtigt. Für ein gasgesättigtes Reservoir aus Tabelle \ref{tab:keff-geology} haben wir \(\Keff \approx 6.1\), woraus sich \(Q \approx Q_0 \times (6.1/4.0)^{0.67} \approx 1.35\,Q_0\) ergibt.
	
	Einsetzen der Zahlenwerte:
	
	\[
	\delta g_{\text{res}} \approx 100\ \text{nm/s}^2 \times 1.35 \times 0.35 \times (50/2000)^2 \approx 100 \times 1.35 \times 0.35 \times 0.000625 \approx 0.03\ \text{nm/s}^2.
	\]
	
	Der erhaltene Wert \(\delta g_{\text{res}} \approx 0.03\ \text{nm/s}^2 = 3 \times 10^{-11}\ \text{m/s}^2\) (d.h. \(3 \times 10^{-12}\ g\)) liegt an der Empfindlichkeitsgrenze moderner supraleitender Gravimeter wie dem iGrav (typische Empfindlichkeit \(10^{-11}\ \text{m/s}^2\), d.h. etwa \(10^{-12}\ g\)). Ein solches Signal könnte durch Stapelung der Daten über mehrere Gezeitenzyklen nachgewiesen werden. Für ein wassergesättigtes Reservoir (\(\Keff \approx 4.5\)) wäre das Signal etwa halb so groß (\( \approx 1.5\times10^{-11}\ \text{m/s}^2\), d.h. \(1.5\times10^{-12}\ g\)), was die Detektion schwieriger macht, aber bei längeren Beobachtungszeiten prinzipiell möglich bleibt. Somit liegt die direkte Registrierung von Gezeitenresonanzen aus tiefen Reservoiren an der Grenze der aktuellen technologischen Fähigkeiten, was das Fehlen zuverlässiger Beobachtungen bis heute erklärt.
	% ========= ENDE DER KORREKTUR =========
	
	\subsection{Beziehung zwischen Resonanzfrequenz und Koordinationseffizienz}
	\label{subsec:frequency_keff}
	
	Gleichung (25) verknüpft den Resonanzgütefaktor \(Q\) mit \(\Keff\), aber für eine vollständige Identifizierung des Reservoirs ist auch die Abhängigkeit der Resonanzfrequenz \(\omega_0\) von \(\Keff\) erforderlich. Im Rahmen des „Feder-Masse-Dämpfer“-Modells, bei dem das Reservoir ein effektives Schwingungssystem darstellt, ist die Steifigkeit \(k\) proportional zum Schermodul \(\mu\), das gemäß Gleichung (\ref{eq:keff-properties}) als \(\mu \propto \Keff^{\xi}\) mit \(\xi \approx 0.4\) skaliert. Dann gilt für die Grundmode-Resonanzfrequenz:
	
	\begin{equation}
		\omega_0 = \sqrt{\frac{k}{M_{\text{eff}}}} \propto \sqrt{\mu} \propto \Keff^{\xi/2} \approx \Keff^{0.2}.
		\label{eq:omega_keff}
	\end{equation}
	
	\subsection{Kalibrierungsbeispiel: Samotlor-Feld}
	
	Als konkretes Beispiel betrachten Sie das Samotlor-Ölfeld in Westsibirien — eines der größten der Welt. Es besteht aus mehreren übereinanderliegenden Reservoiren in Kreidesandsteinen in Tiefen von 1,5–2,5 km, mit Porositäten von 15–25\% und Ölsättigungen von 60–80\%. Unter Veröffentlichung von Daten können wir das erwartete $\Keff$ des ölgesättigten Sandsteins abschätzen: Aus Tabelle \ref{tab:keff-geology} ist die Sandsteinmatrix $\Keff \approx 3.8$, Öl $\Keff \approx 5.2$, daher ist die effektive $\Keff$ des Reservoirs (durch Volumenmittelung) \(\approx 0.7 \times 3.8 + 0.3 \times 5.2 \approx 4.2\) (bei einer Porosität von 30\%). Das vorhergesagte $Q$ wäre \(\approx Q_0 \times 4.2^{0.67} \approx 2.8 Q_0\). Wenn das Hintergrundgestein $\Keff \approx 4.0$ hat, ist der Kontrast bescheiden. Das Vorhandensein von Gaskappen (höheres $\Keff$) oder Wasserfüßen (niedrigeres $\Keff$) würde jedoch nachweisbare spektrale Unterschiede erzeugen.
	
	Eine Pilotvermessung über Samotlor mit 20 supraleitenden Gravimetern über ein Jahr könnte testen, ob Gezeitenresonanzen mit bekannten Reservoirgrenzen korrelieren. Ein Erfolg würde einen starken Proof-of-Concept liefern.
	
	\section{Wirtschaftliche Tragfähigkeit und Risikominderung}
	\label{sec:economics}
	
	\subsection{Kostenvergleich mit konventionellen Methoden}
	
	\begin{table}[htbp]
		\centering
		\caption{Wirtschaftlicher Vergleich von Explorationsmethoden}
		\label{tab:economics}
		\small
		\begin{tabular}{lccc}
			\toprule
			\textbf{Methoden} & \textbf{Kosten pro Vermessung} & \textbf{Eindringtiefe} & \textbf{Umweltauswirkung} \\
			\midrule
			3D-Seismik (Land) & \$5–20 Mio. & 5–8 km & Hoch (Vibroseis, Bohren) \\
			3D-Seismik (Offshore) & \$20–100 Mio. & 8–10 km & Hoch (Luftgewehre, Streamer) \\
			Gravimetrie/Magnetik & \$1–3 Mio. & Integriert & Niedrig \\
			\textbf{Gezeitentomographie} & \textbf{\$1–2 Mio.} & \textbf{Volltiefe} & \textbf{Keine (passiv)} \\
			\bottomrule
		\end{tabular}
	\end{table}
	
	Eine regionale Gezeitentomographie-Untersuchung erfordert den Einsatz von 10–20 hochpräzisen Gravimetern/Seismometern (Stationenkosten ca. \$50.000 pro Stück) sowie Datenverarbeitung und Interpretation, insgesamt etwa \$1–2 Millionen. Dies ist zwei Größenordnungen niedriger als Offshore-Seismikvermessungen, die vergleichbare Gebiete abdecken.
	
	\subsection{Risikominderung bei Explorationsbohrungen}
	
	Branchenstatistiken zeigen, dass selbst in reifen Becken 30–50\% der Explorationsbohrungen trocken (nicht kommerziell) sind. Zusätzliche geophysikalische Informationen können dieses Risiko verringern. Basierend auf Analogien aus verbesserter seismischer Bildgebung ist mit einer Reduzierung der Trockenlochwahrscheinlichkeit um 20–30\% durch die Gezeitentomographie erreichbar.
	
	Bei Kosten für tiefe Explorationsbohrungen von \$10–30 Millionen würde das Einsparen von nur 2–3 Trockenlöchern die Vermessungskosten um das 10- bis 45-fache amortisieren.
	
	\subsection{Umwelt- und soziale Vorteile}
	
	Die Gezeitentomographie ist völlig passiv, benötigt keine Explosionsquellen, Vibratoren oder Bohrungen. Sie ist umweltfreundlich und kann in sensiblen Gebieten (Wildnis, Stadtgebiete, Meeresschutzgebiete) eingesetzt werden, in denen aktive Methoden verboten sind.
	
	\section{Wirtschaftliche und gesellschaftliche Auswirkungen der strukturellen Gezeitenüberwachung}
	\label{sec:structural-economics}
	
	\subsection{Kostenvergleich mit konventionellen Methoden der Bauwerksüberwachung}
	
	\begin{table}[htbp]
		\centering
		\caption{Kostenvergleich von Methoden zur Zustandsüberwachung von Bauwerken}
		\label{tab:structural-costs}
		\small
		\begin{tabular}{lccc}
			\toprule
			\textbf{Methoden} & \textbf{Kosten pro Bauwerk} & \textbf{Abdeckung} & \textbf{Kontinuierlich?} \\
			\midrule
			Sichtprüfung & \$10.000–50.000/Jahr & Nur Oberfläche & Nein \\
			Ultraschallprüfung & \$50.000–200.000 & Lokale Punkte & Nein \\
			Faseroptische Dehnungssensoren & \$100.000–500.000 & Vorinstallierte Leitungen & Ja \\
			Beschleunigungsmesser-Netzwerk & \$50.000–150.000 & Diskrete Punkte & Ja \\
			\textbf{Gezeitenüberwachung} & \textbf{\$50.000–100.000 (einmalig)} & \textbf{Vollvolumen} & \textbf{Ja (passiv)} \\
			\bottomrule
		\end{tabular}
	\end{table}
	
	Ein Gezeitenüberwachungssystem für ein großes Bauwerk benötigt 1–3 hochempfindliche Sensoren (Kosten pro Stück ca. \$30.000–50.000), Datenerfassung und Analysesoftware — eine einmalige Investition von \$50.000–100.000. Dies ist \textbf{Größenordnungen geringer} als der Wert des Bauwerks (ein großer Wolkenkratzer kostet \$500 Millionen–\$1 Milliarde) und die potenziellen Kosten eines Versagens.
	
	\subsection{Risikominderung und vermiedene Verluste}
	
	\begin{itemize}
		\item \textbf{Einsturz der Genua-Brücke (2018):} 43 Tote, wirtschaftlicher Verlust auf \$5 Milliarden geschätzt \cite{Genoa2018}.
		\item \textbf{Einsturz des Surfside-Apartments (2021):} 98 Tote, Verluste über \$1 Milliarde \cite{Surfside2021}.
		\item \textbf{Durchschnittliche jährliche Kosten von Brückeneinstürzen in den USA:} \$1–2 Milliarden \cite{ASCE2021}.
	\end{itemize}
	
	Eine Verringerung des Ausfallrisikos um 10\% durch Früherkennung würde jährlich Hunderte von Leben und Milliarden von Dollar retten. Die Gezeitentomographie ist mit ihrer kontinuierlichen, nicht-invasiven Überwachung einzigartig positioniert, um solche Frühwarnungen zu liefern.
	
	\subsection{Die Ökonomie der Skalierung: Ein planetares Überwachungsnetzwerk}
	
	Betrachten Sie ein globales Netzwerk, das 10.000 kritische Bauwerke (Wolkenkratzer, Brücken, Dämme, kerntechnische Anlagen) überwacht:
	
	\begin{itemize}
		\item Sensoren: 10.000 × \$50.000 = \$500 Millionen.
		\item Installation und Kalibrierung: \$100 Millionen.
		\item Zentrale Datenverarbeitung und KI-Analyse: \$100 Millionen.
		\item \textbf{Gesamt: \$700 Millionen — weniger als die Kosten eines einzigen Flugzeugträgers.}
	\end{itemize}
	
	Dieses Netzwerk würde:
	\begin{itemize}
		\item Echtzeit-Gesundheitsdaten zu den wichtigsten Bauwerken der Welt liefern.
		\item Vorausschauende Wartung ermöglichen und Milliarden an Reparaturkosten einsparen.
		\item Sofortige Bewertung nach Katastrophen ermöglichen, ohne Inspektoren in Gefahrenzonen zu schicken.
		\item Ein historisches Archiv über das Verhalten von Bauwerken unter Gravitationsbelastung erstellen.
	\end{itemize}
	
	Die gesellschaftliche Kapitalrendite würde sich über Jahrzehnte in Billionen von Dollar und Tausenden geretteten Leben messen.
	
	\section{Experimentelle Roadmap}
	\label{sec:roadmap}
	
	\subsection{Phase 1: Koordinationsmineral-Datenbank (1–2 Jahre)}
	
	\begin{itemize}
		\item Zusammenstellung von Labormessungen der elastischen Moduln, Dichte und Porosität für repräsentative Gesteinstypen und Fluide.
		\item Kalibrierung der $\Keff$-Werte unter Verwendung von Skalierungsbeziehungen und dem universellen Fehlergesetz.
		\item Befüllung des koordinationsbasierten Mineral-Klassifikationssystems (Tabelle \ref{tab:keff-geology} erweitert auf Hunderte von Einträgen).
	\end{itemize}
	
	\subsection{Phase 2: Pilot-Feldstudie (2–3 Jahre)}
	
	\begin{itemize}
		\item Auswahl einer gut charakterisierten Kohlenwasserstoffprovinz mit vorhandenen seismischen und Bohrlochdaten (z.B. Westsibirien, Permian Basin, Nordsee).
		\item Einsatz eines Netzwerks von 15–20 supraleitenden Gravimetern (z.B. iGrav-Typ, Empfindlichkeit \(10^{-11}\ \text{m/s}^2\), d.h. etwa \(10^{-12}\ g\)) und Breitbandseismometern über einem 100×100 km Gebiet (Stationenabstand ~25–30 km).
		\item Kontinuierliche Aufzeichnung über mindestens ein Jahr, um mehrere Gezeitenzyklen zu erfassen und Rauschen durch Stapelung zu unterdrücken.
		\item Extraktion der Gezeitenharmonischen unter Verwendung der Methode der kleinsten Quadrate und bekannter Ephemeriden. Berücksichtigung der ozeanischen Gezeitenbelastung mit globalen Modellen (z.B. TPXO) und Subtraktion des modellierten Effekts.
		\item Berechnung der Residuenspektren und Identifizierung von Resonanzpeaks.
		\item Vergleich mit bekannten Reservoirpositionen zur Validierung der Methode.
	\end{itemize}
	
	\subsubsection{Rauschquellen und Minderung}
	
	Potenzielle Rauschquellen umfassen:
	\begin{itemize}
		\item Atmosphärische Druckschwankungen — können durch Luftdruckkorrektur und lokale Drucksensoren reduziert werden.
		\item Mikroseismisches Rauschen (Ozeanwellen, kulturelle Aktivitäten) — wird durch Auswahl ruhiger Standorte und Verwendung von Array-Verarbeitung reduziert.
		\item Instrumentelle Drift — durch Kalibrierung und Vergleich mit Referenzstationen korrigiert.
		\item Hydrologische Belastung (Grundwasser, Schnee) — kann mit lokalen Wetterdaten modelliert werden.
	\end{itemize}
	
	Moderne Datenverarbeitungstechniken (Wavelet-Entrauschung, mehrkanalige Singulärspektrumanalyse) können das Signal-Rausch-Verhältnis weiter verbessern.
	
	\subsection{Phase 3: Entwicklung von Inversionsalgorithmen (3–4 Jahre)}
	
	\begin{itemize}
		\item Entwicklung von 3D-tomographischen Inversionscodes, die die $\Keff$-Parametrisierung einbeziehen.
		\item Validierung an synthetischen Daten und Verfeinerung der Regularisierungsstrategien.
		\item Integration mit komplementären Daten (Magnetotellurik, passive Seismik) für gemeinsame Inversion.
	\end{itemize}
	
	\subsection{Phase 4: Kommerzielle Bereitstellung (4–6 Jahre)}
	
	\begin{itemize}
		\item Aufbau eines kommerziellen Dienstes, der regionale Gezeitentomographie-Untersuchungen anbietet.
		\item Integration in Explorationsabläufe neben seismischen Methoden.
		\item Ausweitung auf Mineralforschung, geothermische Ressourcenbewertung und CO$_2$-Sequestrierungsüberwachung.
	\end{itemize}
	
	\subsection{Parallele Roadmap: Zustandsüberwachung von Bauwerken}
	
	\subsubsection{Phase S1: Konzeptnachweis an instrumentierten Gebäuden (2–3 Jahre)}
	
	\begin{itemize}
		\item Auswahl von 3–5 hohen Gebäuden mit bestehenden Sensornetzwerken (z.B. in Tokio, San Francisco, Dubai).
		\item Installation von zusätzlichen hochempfindlichen Gravimetern/Beschleunigungsmessern.
		\item Aufzeichnung der Gezeitenantworten über 3–6 Monate.
		\item Extraktion der Modalparameter und Korrelation mit bekannten Struktureigenschaften.
		\item Validierung gegen Finite-Elemente-Modelle.
	\end{itemize}
	
	\subsubsection{Phase S2: Basislinien-Establishment und Datenbank (3–5 Jahre)}
	
	\begin{itemize}
		\item Erweiterung auf 50–100 Bauwerke verschiedener Typen (Wolkenkratzer, Brücken, Dämme).
		\item Erstellung einer „gravitativen Fingerabdruck“-Datenbank für jedes Bauwerk.
		\item Entwicklung automatisierter Verarbeitungspipelines und Anomalieerkennungsalgorithmen.
		\item Kalibrierung von \(K_{\mathrm{eff}}\)-Schwellenwerten für verschiedene Schadenszustände.
	\end{itemize}
	
	\subsubsection{Phase S3: Kontinuierliches Überwachungsnetzwerk (5–10 Jahre)}
	
	\begin{itemize}
		\item Installation permanenter Sensoren an 1.000+ kritischen Bauwerken.
		\item Integration in nationale Frühwarnsysteme.
		\item Bereitstellung von Echtzeit-Gesundheitsdashboards für Ingenieure und Behörden.
		\item Festlegung internationaler Standards für die gezeitenbasierte Bauwerksüberwachung.
	\end{itemize}
	
	\subsubsection{Phase S4: Planetarer Maßstab (10–20 Jahre)}
	
	\begin{itemize}
		\item Ausweitung auf 10.000+ Bauwerke weltweit.
		\item Schaffung eines einheitlichen „erdumspannenden“ Überwachungsnetzwerks.
		\item Nutzung der gesammelten Daten zur Verbesserung von Bauvorschriften und Entwurfspraktiken.
		\item Ermöglichung von grenzüberschreitendem Datenaustausch für die Katastrophenhilfe.
	\end{itemize}
	
	Die Kosten dieses gesamten Programms sind ein Bruchteil der jährlichen Verluste durch Bauwerksversagen. Die Technologie ist bereit; es fehlt nur der Wille, sie umzusetzen.
	
	% =============================================================================
	% DETAILLIERTE PILOTPROJEKTE UND KOMMERZIALISIERUNGSAUSSICHTEN
	% =============================================================================
	\section{Detaillierte Pilotprojekte und Weg zur Kommerzialisierung}
	
	\subsection{Pilotprojekt 1: Gravitationstomographie eines Hochhauses}
	
	\textbf{Konzept:} Analog zur medizinischen Ultraschalluntersuchung behandeln wir das Gebäude als einen „Körper“, die Mondgezeiten als externe „Ultraschallquelle“ und ein Netzwerk hochpräziser Gravimeter/Seismometer als „Sensoren“, die die durch innere Inhomogenitäten (Hohlräume, Risse, verschiedene Materialien) verursachten Verzerrungen der Gezeitenwelle aufzeichnen. Durch Analyse dieser Verzerrungen kann eine 3D-Karte der effektiven gravitativen Steifigkeit (verwandt mit $\Keff$) rekonstruiert werden.
	
	\textbf{Instrumentierung (verbesserte Genauigkeit):}
	\begin{itemize}
		\item 9 supraleitende Gravimeter (z.B. iGrav), platziert im Keller, bei 1/4, 1/2, 3/4 der Höhe und auf dem Dach.
		\item 20 Breitbandseismometer (z.B. Trillium 120), verteilt um den Umfang und an Referenzpunkten.
		\item GPS-Synchronisation und meteorologische Station.
	\end{itemize}
	
	\textbf{Messprogramm:}
	\begin{enumerate}
		\item Kontinuierliche Aufzeichnung über 3 Monate (ein vollständiger Mondzyklus mit Reserve).
		\item Abtastrate 100 Hz, um mögliche Resonanzen zu erfassen.
		\item Kalibrierungszeitraum (2 Wochen) zur Aufzeichnung des Hintergrundrauschens.
		\item Gesteuerte Anregungen (Aufzugsbewegungen, Türschlagen) für Testsignale.
	\end{enumerate}
	
	\textbf{Geschätzte Kosten (detailliert):}
	\begin{table}[htbp]
		\centering
		\caption{Budget für Pilotprojekt 1 (USD)}
		\label{tab:budget-building}
		\small
		\begin{tabular}{lrr}
			\toprule
			\textbf{Posten} & \textbf{Berechnung} & \textbf{Kosten} \\
			\midrule
			\multicolumn{3}{l}{\textbf{Gerätemiete}} \\
			Gravimeter (9 Stk.) & $9\times90\times300$ & 243.000 \\
			Seismometer (20 Stk.) & $20\times90\times105$ & 189.000 \\
			Datenlogger+GPS (29 Stk.) & $29\times90\times50$ & 130.500 \\
			Wetterstation (Kauf) & 1 Stk. & 6.000 \\
			Kabel, Batterien, Ersatzteile & – & 15.000 \\
			\multicolumn{3}{l}{\textbf{Transport \& Logistik}} \\
			Hin- und Rücktransport der Geräte & $2\times7.000$ & 14.000 \\
			Lokaler Transporter Miete & $3\times1.500$ & 4.500 \\
			\multicolumn{3}{l}{\textbf{Installation \& Demontage}} \\
			Rigger/Elektriker (10 Pers$\times$5 Tage$\times$400) & – & 20.000 \\
			Inbetriebnahme (5 Tage$\times$2 Ing.$\times$500) & – & 5.000 \\
			\multicolumn{3}{l}{\textbf{Personal (Gehälter)}} \\
			Projektleiter (3 Monate) & $3\times18.000$ & 54.000 \\
			Obergeophysiker (3 Monate) & $3\times12.000$ & 36.000 \\
			Geophysiker (2 Pers$\times$3 Monate$\times$8.000) & – & 48.000 \\
			Techniker (3 Pers$\times$3 Monate$\times$5.500) & – & 49.500 \\
			Laborant (3 Monate) & $3\times4.000$ & 12.000 \\
			\multicolumn{3}{l}{\textbf{Unterkunft \& Verpflegung}} \\
			Miete für nicht-ortansässige Mitarbeiter (3 Pers$\times$3 Monate$\times$1.200) & – & 10.800 \\
			Tagegeld (6 Pers$\times$90 Tage$\times$30) & – & 16.200 \\
			Krankenversicherung (6 Pers$\times$300) & – & 1.800 \\
			\multicolumn{3}{l}{\textbf{Datenverarbeitung \& Software}} \\
			Speziallizenzen (MATLAB, Signalverarbeitung) & – & 12.000 \\
			Cloud Computing & – & 5.000 \\
			Entwicklung von Inversionsalgorithmen (Auftrag) & – & 20.000 \\
			\multicolumn{3}{l}{\textbf{Sonstiges \& Eventualitäten}} \\
			Verwaltungsgemeinkosten & – & 8.000 \\
			Eventualitäten (15\%) & – & 135.000 \\
			\midrule
			\textbf{Gesamt} & & \textbf{1.080.000} \\
			\bottomrule
		\end{tabular}
	\end{table}
	
	\textbf{Erwartetes Ergebnis:} Erste 3D-Karte der inneren Struktur eines Gebäudes, die allein aus passiven Gezeitenmessungen gewonnen wurde. Identifizierung versteckter Defekte und resonanter Moden. Kommerzieller Wert für die Zustandsüberwachung von Wolkenkratzern, Brücken und Dämmen.
	
	\subsection{Pilotprojekt 2: Regionale Tomographie in einem Karstgebiet}
	
	\textbf{Konzept:} Einsatz eines dichten Netzwerks von Gravimetern und Seismometern über einem gut untersuchten Karstgebiet (z.B. Kungur-Eishöhle, Russland), wo bekannte Hohlräume ein Kalibrierungsziel darstellen. Die Mondgezeiten beleuchten den Untergrund; durch Analyse von Amplituden-/Phasenanomalien und Resonanzfrequenzen wird ein 3D-tomographisches Modell des Karstsystems rekonstruiert.
	
	\textbf{Instrumentierung (verbesserte Genauigkeit):}
	\begin{itemize}
		\item 14 supraleitende Gravimeter.
		\item 35 Breitbandseismometer.
		\item 5 Sensoren, die in der Höhle selbst platziert werden.
		\item Referenzstation 50 km entfernt.
	\end{itemize}
	
	\textbf{Messprogramm:}
	\begin{enumerate}
		\item 6 Monate kontinuierliche Aufzeichnung.
		\item Dichtes Gitter (100$\times$100 km, Abstand 5–10 km).
		\item Synchronisation über GPS.
		\item Parallele Mikroklimaüberwachung zur Druckkorrektur.
	\end{enumerate}
	
	\textbf{Geschätzte Kosten (detailliert):}
	\begin{table}[htbp]
		\centering
		\caption{Budget für Pilotprojekt 2 (USD)}
		\label{tab:budget-karst}
		\small
		\begin{tabular}{lrr}
			\toprule
			\textbf{Posten} & \textbf{Berechnung} & \textbf{Kosten} \\
			\midrule
			\multicolumn{3}{l}{\textbf{Gerätemiete}} \\
			Gravimeter (14 Stk.) & $14\times180\times300$ & 756.000 \\
			Seismometer (35 Stk.) & $35\times180\times105$ & 661.500 \\
			Datenlogger+GPS (49 Stk.) & $49\times180\times50$ & 441.000 \\
			Wetterstationen (2 Stk., Kauf) & – & 12.000 \\
			Bohranlagen für Sensorinstallation (2 Einheiten$\times$6 Monate$\times$5.000) & – & 60.000 \\
			Höhlenausrüstung (Satz) & – & 25.000 \\
			\multicolumn{3}{l}{\textbf{Transport \& Logistik}} \\
			Geländewagen (2 Einheiten$\times$6 Monate$\times$4.000) & – & 48.000 \\
			ATVs (2 Einheiten$\times$6 Monate$\times$1.500) & – & 18.000 \\
			Treibstoff und Ersatzteile & – & 30.000 \\
			Gerätetransport & – & 20.000 \\
			\multicolumn{3}{l}{\textbf{Feldlager \& Verpflegung}} \\
			Zelte, Generatoren, Kühlschränke (Kauf) & – & 40.000 \\
			Essen \& Wasser (15 Pers$\times$180 Tage$\times$25) & – & 67.500 \\
			Koch \& Hilfspersonal (2 Pers$\times$6 Monate$\times$3.500) & – & 42.000 \\
			Satelliteninternet (6 Monate$\times$1.000) & – & 6.000 \\
			Medizinische Unterstützung (Sanitäter, Erste Hilfe) & – & 15.000 \\
			\multicolumn{3}{l}{\textbf{Personal (Gehälter)}} \\
			Expeditionsleiter (6 Monate$\times$20.000) & – & 120.000 \\
			Stellv. Cheftgeophysiker (6 Monate$\times$15.000) & – & 90.000 \\
			Geophysiker (4 Pers$\times$6 Monate$\times$12.000) & – & 288.000 \\
			Techniker (4 Pers$\times$6 Monate$\times$7.000) & – & 168.000 \\
			Höhlenführer (2 Pers$\times$3 Monate$\times$8.000) & – & 48.000 \\
			Fahrermechaniker (2 Pers$\times$6 Monate$\times$5.500) & – & 66.000 \\
			\multicolumn{3}{l}{\textbf{Verarbeitung \& Interpretation}} \\
			Supercomputerzeit (10.000 Knoten‑h$\times$0.15) & – & 1.500 \\
			Lizenzen für Tomographie/Inversionssoftware & – & 40.000 \\
			3D-Visualisierung & – & 20.000 \\
			Expertenberatung & – & 30.000 \\
			\multicolumn{3}{l}{\textbf{Sonstiges \& Eventualitäten}} \\
			Genehmigungen \& Bewilligungen & – & 10.000 \\
			Versicherung (Team + Ausrüstung) & – & 25.000 \\
			Eventualitäten (20\%) & – & 650.000 \\
			\midrule
			\textbf{Gesamt} & & \textbf{3.900.000} \\
			\bottomrule
		\end{tabular}
	\end{table}
	
	\textbf{Erwartetes Ergebnis:} Erstes 3D-tomographisches Bild eines Karstsystems, das allein aus passiven Gezeitendaten abgeleitet wurde. Direkte Verifikation an der bekannten Höhlengeometrie. Kalibrierung der Methode für die Kohlenwasserstoffexploration und Verwerfungskartierung.
	
	\subsection{Skalierung und Kostensenkungsaussichten}
	
	Nach erfolgreichen Pilotprojekten tritt die Technologie in die kommerzielle Phase ein. Wichtige Faktoren, die die Kostensenkung vorantreiben:
	
	\begin{itemize}
		\item \textbf{Fertige Softwareplattform:} Die entwickelten Inversionsalgorithmen werden zu einem wiederverwendbaren Produkt, das die F\&E-Kosten eliminiert (30–40\% Einsparung).
		\item \textbf{Mengenrabatte bei Geräten:} Massenproduktion oder Langzeitmiete reduziert die Stückkosten um 20–30\%.
		\item \textbf{Erfahrene Teams \& standardisierte Verfahren:} Die Feldinstallationszeit verkürzt sich, weniger hochbezahlte Spezialisten werden benötigt.
		\item \textbf{Automatisierte Datenverarbeitung:} Eine Pipeline von Rohdaten zum 3D-Modell reduziert die Interpretationskosten um das 2- bis 3-fache.
	\end{itemize}
	
	\textbf{Projizierte Kosten nach Skalierung:}
	
	\begin{table}[htbp]
		\centering
		\caption{Erwartete Kostensenkung für kommerzielle Projekte}
		\label{tab:scaling}
		\begin{tabular}{lccc}
			\toprule
			\textbf{Projekttyp} & \textbf{Pilotkosten} & \textbf{Kommerzielle Kosten} & \textbf{Reduktionsfaktor} \\
			\midrule
			Gebäudeüberwachung & 1,08 Mio.\$ & 0,3–0,5 Mio.\$ & 2–3× \\
			Regionale Tomographie (100×100 km) & 3,9 Mio.\$ & 1,5–2,0 Mio.\$ & 2–2,5× \\
			\bottomrule
		\end{tabular}
	\end{table}
	
	\textbf{Weitere Anwendungen:} Brücken, Dämme, Tunnel; Kohlenwasserstoff- und Erzexploration; geothermische Reservoirbildgebung; CO$_2$-Speicherungsüberwachung; Verwerfungsüberwachung in seismischen Zonen. Mit einem globalen Netzwerk von 50–100 permanenten Stationen könnten die jährlichen Überwachungskosten pro Standort auf \$50–100k sinken, was die Methode mit konventionellen seismischen Untersuchungen und Bohrungen wettbewerbsfähig macht.
	
	Somit legt die anfängliche Investition in Pilotprojekte den Grundstein für einen profitablen kommerziellen Dienst, der die Sicherheit der Infrastruktur, die Ressourcenexploration und die geodynamische Überwachung adressiert.
	
	\section{Fazit und Ausblick}
	\label{sec:conclusion}
	
	Dieser Anhang hat gezeigt, dass Gezeitenwellen ein natürliches, passives und global verfügbares Sondierungssignal für die Erforschung von Planeteninneren darstellen. Im YUCT-Rahmen wird die Gezeitenantwort direkt mit der unterirdischen Verteilung der Koordinationseffizienz $\Keff$ verknüpft, die als universeller Indikator für geologische Eigenschaften dient.
	
	Unabhängige Forschungsarbeiten — einschließlich direkter Beobachtungen von Gezeitenresonanzen \cite{Sibgatulin2014}, Anwendungen der Gezeitentomographie in der Planetenforschung \cite{RoviraNavarro2025}, unterirdischer Gezeitenmessungen \cite{Rogister2024} und planetarer Gezeitenstudien \cite{Briaud2025} — liefern starke empirische Unterstützung für die Grundprinzipien dieses Ansatzes.
	
	Wichtige Innovationen der auf YUCT basierenden Methode umfassen:
	
	\begin{enumerate}
		\item Einen einheitlichen theoretischen Rahmen, der die Gezeitenantwort über das universelle Fehlergesetz $\eps = \kappac \alpha \Keff^{-\beta}$ mit $\Keff$ verknüpft.
		\item Die explizite Einbeziehung der YUCT-Sektorkopplung (Sektoren 0, 34, 1) über den Lagrange-Term $\kappa_{0,34} \mathrm{Tr}(\Psi_{0,34} \cdot O_0 \cdot O_{34}^\dagger)$.
		\item Ein koordinationsbasiertes Mineral-Klassifikationssystem, das die direkte Identifizierung von Fluidtypen aus Resonanzsignaturen ermöglicht.
		\item Nachgewiesene wirtschaftliche Tragfähigkeit durch Risikominderung bei Explorationsbohrungen.
		\item Erweiterung auf die Zustandsüberwachung von Bauwerken, die eine kontinuierliche, passive, nicht-invasive Bewertung von Wolkenkratzern, Brücken und Dämmen ermöglicht.
		\item Einen Rahmen für die Erdbebenvorhersage basierend auf der Überwachung von \(K_{\mathrm{eff}}\) in Verwerfungszonen.
		\item Das Potenzial für Tsunamivorhersage durch Netzwerke von Meeresboden-Gravimetern.
	\end{enumerate}
	
	\subsection{Planetare Anwendungen}
	
	Über die Erde hinaus ist die Methodik direkt auf andere Himmelskörper mit verfügbarer Gezeitenkraft anwendbar:
	
	\begin{itemize}
		\item \textbf{Mond:} Die Gezeitendeformation durch die Erde könnte die innere Struktur des Mondes untersuchen.
		\item \textbf{Mars:} Sonnengezeiten und durch Phobos induzierte Gezeiten könnten Variationen der Krustendicke einschränken.
		\item \textbf{Eismonde (Europa, Ganymed, Enceladus):} Gezeitentomographie, wie von \cite{RoviraNavarro2025} vorgeschlagen, könnte unterirdische Ozeane und Variationen der Eisschilddicke abbilden.
		\item \textbf{Exoplaneten:} Zukünftige Beobachtungen der Gezeitendeformation von Exoplaneten könnten die innere Struktur einschränken.
	\end{itemize}
	
	Der YUCT-Rahmen bietet eine einheitliche Sprache für die Interpretation von Gezeitenantworten über all diese Himmelskörper hinweg und verwandelt Gezeitenbeobachtungen von Rauschen in Signal und von Signal in quantitative tomographische Bilder.
	
	Über seine geophysikalischen Anwendungen hinaus eröffnet die Gezeitentomographie eine neue Grenze im \textbf{Bauingenieurwesen und der Infrastruktursicherheit}. Zum ersten Mal haben wir eine Methode, um kontinuierlich, passiv und nicht-invasiv den strukturellen Gesundheitszustand der größten und kritischsten Bauwerke der Menschheit zu überwachen — von den höchsten Wolkenkratzern bis zu den längsten Brücken, von antiken Monumenten bis zu nuklearen Containments. Dieselben Gravitationswellen, die Planeten seit Milliarden von Jahren formen, können nun genutzt werden, um die Werke der menschlichen Zivilisation zu schützen. Dies ist nicht nur eine inkrementelle Verbesserung; es ist ein Paradigmenwechsel in der Art und Weise, wie wir die gebaute Umwelt verstehen, erhalten und schützen.
	
	\section*{Danksagungen}
	
	Der Autor dankt den Teilnehmern des YUCT-Seminars für anregende Diskussionen und würdigt die bahnbrechende Arbeit von Sibgatulin, Rovira-Navarro, Rogister, Tanaka, Tolstoy und ihren Kollegen, deren unabhängige Untersuchungen wesentliche empirische Grundlagen für den hier entwickelten Ansatz gelegt haben.
	
	\section*{Datenverfügbarkeit}
	
	Alle Berechnungen, Codes und zusätzlichen Materialien sind verfügbar unter \url{https://github.com/Alexey-Yakushev-YUCT/YPSDC}.
	
	\begin{thebibliography}{99}
		
		\bibitem{Yakushev2026a}
		A. V. Yakushev, \emph{Yakushev Unified Coordination Theory (YUCT)}, Zenodo, \url{https://doi.org/10.5281/zenodo.18444599} (2026).
		
		\bibitem{AppendixA}
		Anhang A: Praktischer Leitfaden zur Verwendung von YUCT V35.0 für interdisziplinäre Modellierung und Vorhersagen.
		
		\bibitem{AppendixC}
		Anhang C: YUCT-Chemie – Vollständige mathematische Grundlagen.
		
		\bibitem{AppendixL}
		Anhang L: Fraktale Koordinationsfehlerskalierung – Ein universelles Gesetz von der DNA bis zur Kosmologie.
		
		\bibitem{AppendixP}
		Anhang P: Temperaturabhängigkeit der Gravitation.
		
		\bibitem{Sibgatulin2014}
		V. G. Sibgatulin, S. A. Peretokin, A. A. Kabanov, ``Resonances of gravitational tides and their effect on geological environment,'' \emph{Earth Science Frontiers} 21(4), 303–310 (2014). doi:10.13745/j.esf.2014.04.030
		
		\bibitem{RoviraNavarro2025}
		M. Rovira-Navarro, I. Matsuyama, D. Dirkx, A. Berne, D. Calliess, S. Fayolle, ``Prospects of Using Tidal Tomography to Constrain Ganymede's Interior,'' \emph{Geophysical Research Letters} 52(11), e2025GL114708 (2025).
		
		\bibitem{Rogister2024}
		Y. Rogister, J. Hinderer, U. Riccardi, S. Rosat, ``Subsurface tidal gravity variation and gravimetric factor,'' \emph{Geophysical Journal International} 238(2), 848–859 (2024).
		
		\bibitem{Briaud2025}
		A. Briaud, A. Stark, H. Hussmann, H. Xiao, J. Oberst, A. Rivoldini, ``New Insights from MESSENGER Data on Mercury's Tidal Response and Internal Structure,'' EPSC-DPS Joint Meeting 2025, Helsinki, Finland (2025).
		
		\bibitem{Tanaka2012}
		S. Tanaka, ``Tidal triggering of earthquakes prior to the 2011 Tohoku-Oki earthquake,'' \emph{Geophysical Research Letters} 39, L00G26 (2012).
		
		\bibitem{Tolstoy2013}
		M. Tolstoy, ``Tidal triggering of microearthquakes on the Juan de Fuca Ridge,'' \emph{Geophysical Research Letters} 40, 2013GL057889 (2013).
		
		\bibitem{RAN2025}
		Russische Akademie der Wissenschaften, ``Lunar distance correlation with earthquake probability,'' \emph{Doklady Earth Sciences} 512, 45–52 (2025).
		
		\bibitem{DeformationPrecursors2024}
		K. Chen et al., ``Step-like volumetric strain changes before moderate earthquakes: Evidence for tidal triggering,'' \emph{Journal of Geophysical Research: Solid Earth} 129, e2024JB028456 (2024).
		
		\bibitem{NOAA2024}
		NOAA, ``Tsunami Warning Systems and DART Buoy Networks,'' \url{https://www.tsunami.noaa.gov/} (2024).
		
		\bibitem{USGS2024}
		USGS, ``Earthquake Hazards Program,'' \url{https://earthquake.usgs.gov/} (2024).
		
		\bibitem{Genoa2018}
		Italienisches Ministerium für Infrastruktur, ``Report on the Genoa Bridge Collapse,'' (2018).
		
		\bibitem{Surfside2021}
		NIST, ``Champlain Towers South Collapse Investigation,'' (2021).
		
		\bibitem{ASCE2021}
		American Society of Civil Engineers, ``2021 Report Card for America's Infrastructure,'' ASCE (2021).
		
		\bibitem{Farrell1972}
		W. E. Farrell, ``Earth tides: a review,'' \emph{Reviews of Geophysics} 10, 761–797 (1972).
		
		\bibitem{Agnew2007}
		D. C. Agnew, ``Earth tides,'' in \emph{Treatise on Geophysics}, Vol. 3, 163–195 (2007).
		
		\bibitem{Wahr1995}
		J. Wahr, ``Earth tides,'' in \emph{Global Earth Physics}, 40–46 (1995).
		
	\end{thebibliography}	
\end{document}